% Tessaris specialist-department product and technical architecture.
% Planned architecture recorded on 24 September 2026.
% Intended for inclusion in the Tessaris technical handover.
% Parent-document packages: array, booktabs, enumitem, hyperref, longtable,
% tabularx, url and xcolor.

\section{Installable Specialist Departments and Business Capability Packs}
\label{sec:tessaris-specialist-departments}

\subsection{Product thesis}

Tessaris is intended to make business capabilities that are normally available
only to larger organisations accessible to small and medium-sized businesses.
The product is not merely a collection of prompts that can draft isolated
documents.  Its advantage is that it already holds an authorised, persistent
understanding of the business: its customers, people, suppliers, contracts,
terms, projects, communications, financial evidence, policies, decisions and
operating history.  When a specialist issue arises, Tessaris can assemble the
relevant context and route it to the appropriate specialist capability without
requiring the owner to reconstruct the matter from the beginning.

The intended proposition is:

\begin{quote}
Enterprise-grade specialist guidance, coordination and operating capacity for
businesses that could not normally maintain every specialist function in-house.
\end{quote}

The names \emph{department} and \emph{Pilot} describe product organisation and
workflow ownership.  They do not represent a claim that Tessaris itself is a
solicitor, accountant, recruiter, safety professional or other regulated
professional.  A specialist department supplies AI-assisted business guidance,
document preparation, analysis and workflow automation.  Outputs can contain
errors and do not replace qualified advice or any human decision legally
required in the applicable jurisdiction.

This section records the intended architecture.  Unless a component is
separately identified in the implementation record, the specialist departments
described here are planned capabilities rather than claims about the currently
delivered application.

\subsection{Capability layers}

The business operating system is divided into three layers:

\begin{enumerate}[leftmargin=*]
  \item \textbf{Core business departments.} Sales, Marketing, Customer Service,
        Finance assistance, Operations, People/HR, Product, Work \& Schedule and
        basic reporting form the standard day-to-day operating layer.
  \item \textbf{Installable specialist departments.} IT, Cybersecurity, Legal,
        Compliance, Procurement, Recruitment, Data, Design, Engineering, Risk,
        Project Management and other expert functions are activated when the
        business needs them.
  \item \textbf{Oversight and assurance.} Strategy, Governance, Internal Audit,
        Quality, Tax coordination, Health \& Safety and Business Continuity
        inspect or constrain work performed elsewhere.
\end{enumerate}

Specialist departments should not all become permanent items in the sticky
header.  A \emph{Specialist Departments} library will allow the owner to inspect,
install, configure, activate, pause and remove capability packs.  Active
specialists can appear in the appropriate workspace or be invoked automatically
by the cross-department router when a qualifying matter occurs.

\subsection{Specialist-department catalogue}

\begin{longtable}{@{}p{0.22\textwidth}p{0.70\textwidth}@{}}
\caption{Planned specialist departments and their principal scope.}
\label{tab:specialist-department-catalogue}\\
\toprule
\textbf{Department} & \textbf{Guidance, generation and coordination scope} \\
\midrule
\endfirsthead
\toprule
\textbf{Department} & \textbf{Guidance, generation and coordination scope} \\
\midrule
\endhead
IT \& Systems & Device and software inventory, licences, identity and access,
backups, help-desk work, systems configuration and technology suppliers. \\
Cybersecurity & Local security scanning, vulnerability and configuration
review, exposed services, credential risks, incident preparation, evidence,
recovery coordination and staff awareness. \\
Legal Support & Contract and terms review, document preparation, disputes,
obligations, intellectual-property records, evidence bundles and escalation to
an appropriately qualified professional. \\
Compliance \& Privacy & Data protection, consent, retention, internal policies,
industry requirements, control evidence, assessments and audit preparation. \\
Risk Management & Risk registers, controls, insurance gaps, operational
incidents, dependencies, scenario analysis and treatment tracking. \\
Procurement & Requirement detection, supplier discovery, requests for quote,
comparison, negotiation preparation, purchasing controls, approvals and renewal
management. \\
Vendor Management & Supplier contracts, service levels, dependencies,
performance, failures, pricing changes, concentration risks and alternatives. \\
Recruitment \& Workforce Planning & Capacity diagnosis, role design, sourcing,
approved advertising, candidate workflow, interview scheduling, channel
experiments, offers and onboarding preparation. \\
Project \& Programme Management & Goals, teams, AI and human assignments,
budgets, milestones, dependencies, risks, delivery evidence, re-planning and
outcome tracking. \\
Product Management & Customer needs, requirements, roadmaps, prioritisation,
release plans and feedback-to-product traceability. \\
Research \& Development & Research questions, experiments, prototypes,
technical evaluation, innovation records and evidence-led recommendations. \\
Software Engineering & Applications, integrations, automation, code review,
testing, deployment preparation and maintenance under software-change controls. \\
Design \& Creative & Brand systems, graphics, product design, user experience,
websites, prototypes, campaign assets and approved design libraries. \\
Data \& Analytics & Data quality, reporting, dashboards, forecasting, business
intelligence, metric definitions and provenance. \\
Strategy & Market and competitor analysis, business plans, growth options,
scenario planning, resource allocation and strategic review. \\
Quality Assurance & Standards, inspections, complaints, defects, root-cause
analysis, corrective actions and evidence of closure. \\
Internal Audit & Independent inspection of whether policies, approvals,
controls, records and provider receipts are functioning as designed. \\
Corporate Governance & Board records, policies, decisions, delegated authority,
conflicts, statutory calendars and company obligations. \\
Tax Coordination & Tax calendars, evidence assembly, calculation support and
handoff to an appropriately qualified adviser; not unreviewed filing or advice. \\
Treasury & Cash forecasting, banking evidence, liquidity, currency exposure,
debt, working capital and controlled payment preparation. \\
Insurance & Coverage inventories, gaps, renewals, claims preparation and
supporting evidence. \\
Facilities \& Workplace & Premises, utilities, maintenance, leases, workplace
services, inspections and supplier coordination. \\
Health \& Safety Support & Risk-assessment preparation, training and equipment
records, incident capture, inspection schedules and escalation; not replacement
of competent-person duties. \\
Learning \& Development & Skills assessment, onboarding, learning plans,
competency evidence and renewal schedules. \\
Sustainability / ESG & Energy, waste, emissions, supplier standards, improvement
plans and evidence-backed reporting. \\
Communications \& PR & Editorial planning, announcements, stakeholder updates,
reputation monitoring, response preparation, crisis communications and results. \\
Partnerships & Channel partners, affiliates, alliances, joint initiatives,
commercial evidence and relationship obligations. \\
Business Continuity & Critical-process maps, backup operations, disaster
recovery, emergency plans, exercises and resilience evidence. \\
\bottomrule
\end{longtable}

The recommended first specialist release comprises Legal Support, IT \&
Cybersecurity, Compliance \& Privacy, Procurement, Recruitment, Project
Management, Design \& Software Engineering, Data \& Analytics, Health \& Safety,
Communications \& PR, Strategy and Business Continuity.  These cover common
capabilities that smaller businesses frequently outsource or operate without.

\subsection{Department-pack contract}

Each installed department is a versioned capability pack rather than a name and
a system prompt.  A pack should declare at least:

\begin{itemize}[leftmargin=*]
  \item a stable department and pack identifier, version and publisher;
  \item its purpose, supported tasks and explicit exclusions;
  \item country, jurisdiction, industry and business-size applicability;
  \item one or more supported hosted or local models and their fallback order;
  \item a specialist harness containing workflows, tools, structured outputs,
        evaluation cases, escalation rules and policy constraints;
  \item the Business Container records it may read and the fields it may write;
  \item connected applications, external professionals and provider scopes;
  \item required approvals, eligible automation and non-removable controls;
  \item retention, confidentiality and local-versus-cloud processing rules;
  \item budget, rate and external-spend limits;
  \item required evidence, action receipts and completion criteria; and
  \item health, evaluation, update, suspension and rollback behaviour.
\end{itemize}

The pack can operate through three execution arrangements:

\begin{description}[leftmargin=0pt,style=nextline]
  \item[Internal Pilot]
  Tessaris performs bounded analysis, generation and coordination using the
  normal business context and shared deterministic tools.

  \item[Specialist model and harness]
  A model selected for a domain is combined with a jurisdictional or industrial
  harness.  The harness, rather than the model name alone, defines admitted
  inputs, deterministic checks, tools, limits, required citations, review gates
  and output schema.

  \item[External professional or provider]
  Tessaris prepares a structured matter and its supporting evidence for a
  solicitor, accountant, recruiter, security consultant or other specialist.
  The external party's response and authority remain distinguishable from an AI
  proposal.
\end{description}

No pack may silently expand its own data access, budget or authority because a
new model has been selected.  Model selection and operational authority are
separate controls.

\subsection{Persistent context and specialist work records}

A specialist should receive only the authorised subset of the shared Business
Container required for a matter.  The system should create a canonical
\emph{specialist work record} with:

\begin{itemize}[leftmargin=*]
  \item matter or project identity, originating event and owning department;
  \item objective, scope, jurisdiction, business entity and responsible owner;
  \item linked customers, people, suppliers, contracts and projects;
  \item immutable references to source messages, documents, photographs,
        transactions, policies and previous decisions;
  \item participating human and AI team members;
  \item budget, deadlines, milestones, dependencies and risk classification;
  \item proposals, reviews, approvals, actions and provider receipts;
  \item unresolved questions, professional escalations and final outcome; and
  \item provenance sufficient to reproduce why each recommendation was made.
\end{itemize}

This record is the primary difference between an embedded specialist and a
stand-alone chat.  A new conversation must not lose the business's approved
terms, earlier decisions or the evidence already attached to the matter.

\subsection{Automatic cross-department routing}

Users should not need to know which specialist to open.  Existing Pilots and
inbound connectors emit typed business events.  A specialist-routing service
evaluates the event against installed pack manifests and proposes or creates a
matter with the minimum necessary context.  Routing does not itself grant
authority to communicate externally or make an irreversible decision.

For example, a customer threat can follow this path:

\begin{enumerate}[leftmargin=*]
  \item Customer Service preserves the complete communication and opens an
        incident rather than continuing an improvised argument.
  \item Legal Support receives the customer, accepted terms, contract,
        communications, invoices, work records, photographs and relevant policy.
  \item Finance identifies disputed invoices and payment evidence.
  \item Operations supplies delivery records, assignments and completion proof.
  \item Legal Support identifies missing facts, prepares options and produces an
        approval-gated response or a bundle for an external professional.
\end{enumerate}

Additional routing examples include:

\begin{itemize}[leftmargin=*]
  \item ``We need two more engineers'' routes to People, Recruitment, Finance
        and Operations for capacity, affordability and sourcing.
  \item A supplier price increase routes to Procurement, Finance and Risk.
  \item A new customer portal routes to Product, Design, Engineering, Legal
        Support, Privacy and Cybersecurity.
  \item A journalist enquiry routes to Communications \& PR, Legal Support and
        the business owner.
  \item A fixed-budget expansion routes to Project Management, Procurement,
        People and Finance.
  \item A suspicious attachment or exposed credential routes locally to
        Cybersecurity and preserves evidence before any remediation.
\end{itemize}

The router should return the matched rules, installed pack version, context
selection and required approvals.  It must not invent a specialist capability
when the relevant pack is absent, unhealthy or outside its jurisdiction.

\subsection{Priority specialist workflows}

\subsubsection{Legal Support}

Legal Support may review or prepare terms and conditions, customer and supplier
contracts, employment and contractor documents, website privacy and cookie
materials, obligation calendars, debt-recovery preparation, intellectual-
property records, dispute correspondence and evidence bundles.  Legal packs
should be jurisdiction- and domain-specific, for example UK small-business law,
Irish company and employment law, Spanish commercial operations, construction,
e-commerce and data protection.

The interface must distinguish information and document preparation from advice
or representation by a qualified professional.  Settlement, admission,
termination, signature, filing and contentious external communication remain
review-controlled.  Where professional review is required, Tessaris should
produce a concise issue summary, chronology, documents, missing facts and
questions rather than an unstructured transcript.

\subsubsection{Local IT and Cybersecurity}

Cybersecurity should be continuously available in the desktop or VM product and
should prefer local processing for logs, files, device metadata and secrets.  A
local model and deterministic scanners may inspect software inventory, missing
updates, vulnerable packages, open ports, exposed services, backup health,
access permissions, suspicious changes, phishing indicators, website or domain
configuration and possible credential disclosure.

Scanning, evidence capture and prioritised guidance may be continuous.
Destructive remediation, account suspension, firewall changes, deletion,
quarantine that interrupts the business and active offensive testing require
explicit authority.  Every finding must identify its source, scanner or model,
time, confidence, affected asset, suggested treatment and verification result.

\subsubsection{Procurement and Vendor Management}

The Procurement Pilot may detect demand from projects, schedules, stock and
budgets; discover suppliers; prepare requests for quote; normalise offers;
compare total cost, delivery, quality, warranty and risk; monitor renewal and
price changes; maintain approved suppliers; and prepare negotiation or purchase
recommendations.  Purchase commitments remain subject to business authority and
budget controls.

Procurement is also a potential revenue line through premium subscriptions,
verified supplier introductions, marketplace services, managed procurement or
shared savings.  Any commission, ranking influence or commercial relationship
must be disclosed.  Paid placement must never be presented as an independent
best-value recommendation.

\subsubsection{Recruitment and Workforce Planning}

Recruitment begins with operational demand rather than a job advert.  For a
request such as ``we need more engineers'', the system should inspect workload,
pipeline, skills, availability and budget; compare hiring, contracting and
rescheduling; prepare the role and budget; publish only approved vacancies;
operate approved organic or paid sourcing channels; manage candidates and
interviews; prepare offers and onboarding; and report channel outcomes.

Candidate ranking must use declared job-related criteria, retain provenance and
support challenge and human review.  Posting, paid media, direct candidate
contact, rejection, selection, compensation and offers require the configured
employment controls.  The pack must accommodate data-protection and
anti-discrimination requirements for the deployment jurisdiction.

\subsubsection{Project and Programme Management}

Projects coordinate both human and AI teams around an approved goal.  Each
project holds the success definition, owner, team, budget, authority limits,
milestones, dependencies, documents, decisions, risks, expected outcome, actual
cost and time, and completion evidence.  The Project Pilot may identify delays,
request updates, prepare re-plans and escalate scope or budget variance.  It
connects to Workflow Canvas, Work \& Schedule, People, Procurement, Finance and
Operations rather than creating a separate unlinked task database.

\subsubsection{Communications and PR}

Communications \& PR may maintain the editorial calendar and brand voice,
identify business events worth communicating, prepare announcements, website
and social content, stakeholder updates, press materials, reactive statements
and crisis plans, monitor reputation and measure outcomes.  Routine content may
operate within an approved campaign policy.  Crisis statements, admissions,
regulated claims and legally or reputationally sensitive publications require
human and, where configured, Legal Support review.

\subsection{Authority and human-oversight levels}

Installation and activation are separate.  Every newly activated department
starts in \emph{Review required}.  The owner can later grant more authority to
eligible action classes, but authority is versioned, scoped and revocable.

\begin{longtable}{@{}p{0.22\textwidth}p{0.70\textwidth}@{}}
\caption{Department operating levels.}
\label{tab:specialist-operating-levels}\\
\toprule
\textbf{Level} & \textbf{Permitted behaviour} \\
\midrule
\endfirsthead
\toprule
\textbf{Level} & \textbf{Permitted behaviour} \\
\midrule
\endhead
Review required & Tessaris researches, analyses and prepares drafts or actions,
but nothing leaves the controlled workspace and no external mutation occurs
without an authorised person's approval of the exact payload. \\
Operate within policy & Tessaris performs defined, reversible actions within an
approved policy, data scope, budget, recipient class and rate limit.  Exceptions
return to review. \\
Autonomous routine work & Low-risk, repeatable actions may run without per-action
approval when logging, limits, monitoring, rollback and suspension are present.
This level is not available for prohibited or permanently review-bound actions. \\
\bottomrule
\end{longtable}

Human oversight must be effective rather than ceremonial.  The reviewer needs
the competence, information, time and authority to reject or stop the action.
The interface must show source evidence, uncertainty, important assumptions,
changes since the last approved version and the exact external effect.

The following action classes remain human-controlled unless a later legal and
safety assessment establishes a lawful, narrower implementation:

\begin{itemize}[leftmargin=*]
  \item legal demands, admissions, settlements, representation and contract
        signature;
  \item statutory or tax filing and ambiguous accounting treatment;
  \item candidate or employee rejection, dismissal and other materially
        consequential employment decisions;
  \item safety-critical decisions or assertions that a workplace or activity is
        safe;
  \item transfers of money, bank-detail changes and new financial beneficiaries;
  \item destructive cybersecurity remediation and offensive testing;
  \item crisis, regulated or legally sensitive public statements; and
  \item actions that are irreversible, exceed budget or lack a verified receipt.
\end{itemize}

A customer setting cannot remove a control required by law or by the product's
minimum safety policy.  Increasing autonomy transfers operational discretion
only within the published boundary; it does not convert an unavailable or
unverified capability into an authorised one.

\subsection{Installation and administration experience}

The Specialist Departments library should provide catalogue, installed and
active views.  Before activation, the owner must be able to inspect:

\begin{itemize}[leftmargin=*]
  \item the department's capabilities and exclusions;
  \item jurisdiction, industry and supported business context;
  \item model, harness, publisher, version and evaluation status;
  \item data requested from each Business Container and its purpose;
  \item available connectors and external-professional handoffs;
  \item actions that require approval and those eligible for automation;
  \item non-removable safety boundaries and known limitations;
  \item professional review recommendations or requirements;
  \item retention, local processing and cloud-processing position;
  \item monthly usage and external-spend limits; and
  \item update, rollback and removal consequences.
\end{itemize}

Activation should create the department workspace, default Review-required
authority profile, matter registry, event subscriptions and health checks.  It
must not silently grant a provider scope or copy all business records.  Removal
disables routing and tools but preserves records required by the configured
retention and audit policy.  An \emph{Activate all} operation applies a visible
bundle manifest and requests the combined permissions; it is not a bypass for
department-specific boundaries.

\subsection{Commercial packaging}

The initial packaging hypothesis is:

\begin{description}[leftmargin=0pt,style=nextline]
  \item[Tessaris Business Core]
  Sales, Marketing, Customer Service, Finance assistance, Operations, People/HR,
  Work \& Schedule and basic reporting are included in the standard business
  operating system.

  \item[Tessaris Complete Business]
  An \emph{Activate all} package makes the supported specialist catalogue
  available under one subscription while still requiring individual permission
  and authority configuration.

  \item[Specialist Packs]
  Individual departments, jurisdictions or industries can be installed as
  optional packs.  Pricing may be a small monthly department fee, a premium
  pack subscription or a usage allowance rather than a separate model licence.
\end{description}

Additional revenue can come from professional-review marketplaces,
recruitment-channel management, managed cybersecurity, project or campaign
budgets, and disclosed procurement introductions or shared savings.  Model and
harness cost should be presented as the capability price rather than forcing
the customer to understand token or model infrastructure.  Pricing remains a
commercial hypothesis until usage, support, insurance and provider costs have
been measured.

\subsection{Trust, claims and regulatory position}

The product should not use an ``entertainment purposes only'' statement for
business capabilities designed to influence real work.  Product descriptions,
limitations and demonstrations must accurately reflect the intended use.
Appropriate wording is that the department provides AI-assisted guidance,
document preparation and workflow automation; outputs may contain errors and do
not replace qualified professional advice or legally required human decisions.

Terms and customer-selected autonomy do not replace product controls,
evaluation, clear instructions, incident handling or compliance with applicable
law.  Recruitment and other uses that materially influence people require
particular assessment.  Solely automated decisions that have legal or similarly
significant effects may engage data-protection restrictions and safeguards.
High-risk uses may require risk management, technical documentation, logging,
accuracy, robustness, cybersecurity, effective human oversight and deployer
controls.  The exact product claims, contractual allocation, insurance and
jurisdictional implementation must be reviewed by appropriately qualified legal
and regulatory professionals before public release.

\subsection{Technical service boundaries}

The planned implementation should maintain the following separations:

\begin{enumerate}[leftmargin=*]
  \item \textbf{Catalogue service:} discovers signed pack manifests without
        granting their requested access.
  \item \textbf{Installation service:} verifies compatibility, publisher,
        version, model availability and evaluations, then records consent.
  \item \textbf{Context broker:} produces a purpose-limited evidence bundle from
        authorised Business Container records.
  \item \textbf{Specialist router:} matches typed events to healthy installed
        capabilities and records the reason for the route.
  \item \textbf{Matter service:} owns the canonical specialist work record and
        cross-department references.
  \item \textbf{Harness runtime:} admits only declared tools and structured
        outputs and runs deterministic checks around model proposals.
  \item \textbf{Authority gateway:} compares each proposed action with actor,
        department, action, budget, recipient and risk policies.
  \item \textbf{Approval service:} binds approval to an exact payload hash and
        expires it when the payload or material context changes.
  \item \textbf{Execution adapters:} perform permitted external actions and
        return provider identifiers and read-back evidence.
  \item \textbf{Receipt and audit service:} preserves sources, proposal,
        reviewer, approval, execution result, model and pack versions.
  \item \textbf{Health and suspension service:} disables routing or execution
        when a pack, model, connector, evaluation or jurisdiction is invalid.
\end{enumerate}

The model proposes; the harness validates and structures; the authority gateway
decides whether execution is admissible; a human approves where required; the
adapter acts; and the receipt proves what occurred.  No user-interface label
such as ``department'', ``active'' or ``autonomous'' may bypass that chain.

\subsection{Acceptance criteria for the first release}

The first specialist-department release is complete only when:

\begin{enumerate}[leftmargin=*]
  \item the library distinguishes available, installed, active, unhealthy and
        unavailable-for-jurisdiction packs;
  \item installation displays and records all requested data, tool, provider,
        budget and retention scopes;
  \item activation defaults to Review required and exposes no hidden external
        write authority;
  \item at least one typed business event opens a canonical specialist matter
        with a minimal, inspectable context bundle;
  \item the user can trace every included fact to its source record;
  \item model unavailability, low confidence or missing evidence causes an
        explicit limitation or escalation rather than a fabricated answer;
  \item exact-payload approvals are invalidated by material changes;
  \item every external action produces a receipt or is reported as unverified;
  \item permanent human-control actions cannot be enabled through ordinary
        customer settings;
  \item pack updates can be evaluated, staged, rolled back and audited;
  \item removal stops new routing and execution without destroying retained
        matter evidence; and
  \item marketing and in-product wording accurately separates assistance,
        automation and professional responsibility.
\end{enumerate}

The resulting system is a coordinated organisation of persistent,
context-aware capabilities operating under the business's data permissions,
authority, budget and evidence controls.  This in-house context, controlled
handoff and durable institutional memory---rather than generic document
generation alone---is the intended competitive advantage.
